\documentclass[12pt, a4paper, twoside]{book}

% ─── Packages ───────────────────────────────────────────────────────────────
\usepackage[a4paper, inner=3.0cm, outer=2.5cm, top=2.8cm, bottom=2.8cm, includeheadfoot, headheight=24pt]{geometry}
\usepackage[T1]{fontenc}
\usepackage[utf8]{inputenc}
\usepackage{amsmath, amsthm}
\usepackage[charter, cal=cmcal]{mathdesign}
\usepackage[scaled=0.92]{helvet}
\usepackage{mathtools}
\usepackage{physics}
\usepackage{graphicx}
\usepackage{xcolor}
\usepackage{tikz}
\usetikzlibrary{arrows.meta, decorations.pathmorphing, shapes.geometric,
                positioning, calc, patterns, backgrounds, decorations.markings}
\usepackage{pgfplots}
\pgfplotsset{compat=1.18}
\usepackage{tcolorbox}
\tcbuselibrary{skins, breakable, theorems}
\usepackage{enumitem}
\usepackage{booktabs}
\usepackage{array}
\usepackage{multirow}
\usepackage{fancyhdr}
\usepackage{titlesec}
\usepackage{hyperref}
\usepackage{wrapfig}
\usepackage{float}
\usepackage{caption}
\usepackage{subcaption}
\usepackage{epigraph}
\usepackage{lettrine}
\usepackage{mdframed}
\usepackage{setspace}
\usepackage{parskip}
\usepackage{siunitx}

% ─── Colours ────────────────────────────────────────────────────────────────
\definecolor{deep}{HTML}{0F172A}      % Deep Slate/Charcoal for body text/titles
\definecolor{accent}{HTML}{E11D48}    % Crimson/Coral for problem counters
\definecolor{teal}{HTML}{0D9488}      % Electric Teal for rules and highlights
\definecolor{gold}{HTML}{D97706}      % Gold for puzzle borders
\definecolor{midblue}{HTML}{334155}   % Muted Blue-Gray for section headings
\definecolor{purple}{HTML}{6D28D9}    % Purple for reflection borders
\definecolor{lightbg}{HTML}{F8FAFC}   % Very light gray-blue for exploration bg
\definecolor{warmbg}{HTML}{FEF3C7}    % Warm gold-yellow for puzzle bg
\definecolor{greenbg}{HTML}{F0FDF4}   % Light green for thought/keyidea bg

% ─── Page geometry fix ──────────────────────────────────────────────────────
\setlength{\headheight}{24pt}

% ─── Hyperref ───────────────────────────────────────────────────────────────
\hypersetup{
  colorlinks=true, linkcolor=midblue, urlcolor=accent, citecolor=teal,
  pdftitle={Discovery Physics: Journey 2 — Motion and Its Causes},
  pdfauthor={Discovery Institute},
}

% ─── Page style ─────────────────────────────────────────────────────────────
\pagestyle{fancy}
\fancyhf{}
\fancyhead[LE]{\small\textcolor{deep}{\textit{Discovery Physics}}}
\fancyhead[RO]{\small\textcolor{deep}{\textit{\leftmark}}}
\fancyfoot[C]{\small\textcolor{deep}{\thepage}}
\renewcommand{\headrulewidth}{0.4pt}
\renewcommand{\headrule}{\hbox to\headwidth{\color{teal}\leaders\hrule height \headrulewidth\hfill}}

% ─── Chapter / Section titles ───────────────────────────────────────────────
\titleformat{\chapter}[display]
  {\normalfont\huge\bfseries\color{deep}}
  {\textcolor{accent}{\chaptertitlename\ \thechapter}}{12pt}{\Huge}
  [\vspace{4pt}\textcolor{teal}{\titlerule[1.5pt]}]

\titleformat{\section}
  {\normalfont\Large\bfseries\color{midblue}}
  {\textcolor{accent}{\thesection}}{1em}{}

\titleformat{\subsection}
  {\normalfont\large\bfseries\color{deep}}
  {\textcolor{teal}{\thesubsection}}{1em}{}

% ─── Custom boxes ───────────────────────────────────────────────────────────
\newtcolorbox{puzzlebox}[1][]{
  enhanced, breakable,
  colback=warmbg, colframe=gold,
  fonttitle=\bfseries\large\color{deep}, title={The Puzzle},
  attach boxed title to top left={yshift=-2mm, xshift=4mm},
  boxed title style={colback=gold, colframe=gold}, #1}

\newtcolorbox{explobox}[1][]{
  enhanced, breakable,
  colback=lightbg, colframe=midblue,
  fonttitle=\bfseries\color{white}, coltitle=white, title={Exploration},
  attach boxed title to top left={yshift=-2mm, xshift=4mm},
  boxed title style={colback=midblue, colframe=midblue}, #1}

\newtcolorbox{discoverybox}[1][]{
  enhanced, breakable,
  colback=white, colframe=accent,
  fonttitle=\bfseries\large\color{white}, coltitle=white, title={Discovery},
  attach boxed title to top left={yshift=-2mm, xshift=4mm},
  boxed title style={colback=accent, colframe=accent}, drop shadow, #1}

\newtcolorbox{namebox}[1][]{
  enhanced,
  colback=deep!5, colframe=deep,
  fonttitle=\bfseries\color{white}, coltitle=white, title={What This Is Called},
  attach boxed title to top left={yshift=-2mm, xshift=4mm},
  boxed title style={colback=deep, colframe=deep}, #1}

\newtcolorbox{reflectionbox}[1][]{
  enhanced, breakable,
  colback=purple!5, colframe=purple,
  fonttitle=\bfseries\color{white}, coltitle=white, title={Reflection}, #1}

\newtcolorbox{historybox}[1][]{
  enhanced, breakable,
  colback=gold!8, colframe=gold!70!brown,
  fonttitle=\bfseries\color{deep}, title={Historical Note}, #1}

% ─── Problem counter ────────────────────────────────────────────────────────
\newcounter{problem}[chapter]
\newcommand{\problem}[1][]{%
  \refstepcounter{problem}%
  \medskip\noindent
  \textcolor{accent}{\textbf{Problem \theproblem.}}\ %
  \ifx&#1&\else\textit{(#1)}\ \fi}

\newcommand{\keyidea}[1]{%
  \begin{tcolorbox}[enhanced, colback=teal!8, colframe=teal,
    left=4mm, fontupper=\bfseries\color{deep}]#1\end{tcolorbox}}

% ─── Theorem environments ───────────────────────────────────────────────────
\newtheoremstyle{disc}{8pt}{8pt}{\itshape}{}{\bfseries\color{midblue}}{.}{.5em}{}
\theoremstyle{disc}
\newtheorem{principle}{Principle}[chapter]
\newtheorem{question}{Question}[chapter]

% ─── Spacing ────────────────────────────────────────────────────────────────
\setlength{\epigraphwidth}{0.65\textwidth}
\renewcommand{\epigraphflush}{center}
\setlength{\parskip}{6pt}
\setlength{\parindent}{0pt}
\onehalfspacing

% ════════════════════════════════════════════════════════════════════════════
\begin{document}


\setcounter{chapter}{1}
\setcounter{section}{0}

\section{Module 2.1 — Inventing the Language of Motion}
\label{sec:language_motion}

\epigraph{
  \itshape To measure is to know.
}{--- Lord Kelvin}

\subsection{The Central Question}

\begin{tcolorbox}[colback=deep!5, colframe=deep, fontupper=\large\bfseries\color{deep}]
A stone falls. A planet orbits. A car accelerates. Can we describe all of these with the same language?
\end{tcolorbox}

\subsection{Position, Distance, and Displacement}

Before we can describe motion, we need a way to describe \textit{where} something is. This requires a reference point and a coordinate system---a way of assigning numbers to positions.

For motion along a line, a single number suffices: the position $x$, measured from some agreed origin. The distance between two positions is $|x_2 - x_1|$. But distance alone is not enough: we also need direction. A displacement of $+5$ metres (moving right) is different from a displacement of $-5$ metres (moving left).

\begin{namebox}
\textbf{Displacement} is the change in position, including direction:
\begin{equation}
  \Delta x = x_{\text{final}} - x_{\text{initial}}
\end{equation}
It is positive if the motion is in the positive direction, negative otherwise. It is a \textbf{vector}---a quantity with both magnitude and direction.
\end{namebox}

\begin{explobox}
\problem A runner starts at position $x = 0$, runs to $x = 100\,\text{m}$, then returns to $x = 40\,\text{m}$.
\begin{enumerate}[label=(\alph*)]
  \item What is the total distance travelled?
  \item What is the displacement from start to finish?
  \item Why are these different, and which is more useful for describing the runner's final state?
\end{enumerate}

\problem A ball is thrown upward and returns to its starting point. What is its displacement? What is the distance it travelled?

\problem Why is it not enough to say ``the car moved 50 metres''? What additional information is needed for a complete description?
\end{explobox}

\subsection{Speed and Velocity}

Distance and displacement describe where something is. Speed and velocity describe how fast it is getting there.

\begin{equation}
  \text{average speed} = \frac{\text{distance travelled}}{\text{time taken}}
  \qquad
  \bar{v} = \frac{\Delta x}{\Delta t}
\end{equation}

But again, direction matters. A car moving at \SI{60}{km/h} northward and a car moving at \SI{60}{km/h} southward have the same speed but opposite velocities. \textbf{Velocity} is speed with a direction.

\begin{explobox}
\problem A train travels 120 km in 90 minutes. What is its average speed in km/h? In m/s?

\problem A person walks 3 km east, then 4 km north. The journey takes 2 hours.
\begin{enumerate}[label=(\alph*)]
  \item What is the total distance walked?
  \item What is the displacement (magnitude and direction)?
  \item What is the average speed?
  \item What is the average velocity (magnitude and direction)?
\end{enumerate}

\problem The position of a particle is given by the data below. Compute the average velocity over each interval, and then over the full time.

\begin{center}
\begin{tabular}{ccc}
\toprule
Time (s) & Position (m) & Avg. velocity (m/s) \\
\midrule
0 & 0 & --- \\
1 & 5 & \\
2 & 18 & \\
3 & 39 & \\
4 & 68 & \\
\bottomrule
\end{tabular}
\end{center}
What pattern do you notice in the average velocities?
\end{explobox}

\subsection{The Problem of Instantaneous Velocity}

Average velocity works over a time interval. But what about \textit{right now}, at this instant? A speedometer reads your speed at this moment, not your average over the last hour. What does ``speed right now'' actually mean?

This is not a trivial question. At a single instant, no time has elapsed and no distance has been covered. The ratio $\Delta x / \Delta t$ gives $0/0$---meaningless.

The resolution came from considering what happens as the time interval $\Delta t$ is made smaller and smaller:

\begin{explobox}
\problem A ball's position is $x(t) = 5t^2$ (in metres, with $t$ in seconds).

\begin{enumerate}[label=(\alph*)]
  \item Compute the average velocity between $t = 2$ and $t = 3$.
  \item Compute the average velocity between $t = 2$ and $t = 2.5$.
  \item Compute the average velocity between $t = 2$ and $t = 2.1$.
  \item Compute the average velocity between $t = 2$ and $t = 2.01$.
  \item What value do these approach as the interval gets smaller?
\end{enumerate}

\problem For the same $x(t) = 5t^2$, try to find a formula for the average velocity between $t = 2$ and $t = 2 + h$, where $h$ is a small positive number. What happens as $h \to 0$?

\problem Repeat for $x(t) = 3t^2 + 2t$. What is the velocity at $t = 1$?
\end{explobox}

\begin{discoverybox}
What you are computing in the limit above is the \textbf{derivative} of position with respect to time. The instantaneous velocity is:
\begin{equation}
  v(t) = \lim_{\Delta t \to 0} \frac{x(t + \Delta t) - x(t)}{\Delta t} = \frac{dx}{dt}
\end{equation}
This is the same limit that defines the slope of a curve at a point. The velocity is the slope of the position-time graph.

For $x = kt^n$, the derivative gives $v = knt^{n-1}$. This is a rule you will use constantly---but notice how it \textit{emerged} from the problem rather than being handed to you.
\end{discoverybox}

\subsection{Acceleration: The Rate of Change of Velocity}

If velocity is the rate of change of position, what is the rate of change of velocity? This quantity---how quickly the velocity itself is changing---is called \textbf{acceleration}.

\begin{equation}
  a(t) = \frac{dv}{dt} = \frac{d^2x}{dt^2}
\end{equation}

\begin{figure}[H]
  \centering
  \includegraphics[width=0.9\textwidth]{figures/fig_j2_kinematics}
  \caption{Left column: constant velocity (position grows linearly, velocity is flat). Right column: uniform acceleration (position grows as $t^2$, velocity grows linearly). The area under the velocity-time graph equals the displacement.}
  \label{fig:kinematics}
\end{figure}

\begin{explobox}
\problem A car's velocity is measured every second:

\begin{center}
\begin{tabular}{cc}
\toprule
Time (s) & Velocity (m/s) \\
\midrule
0 & 0 \\
1 & 3 \\
2 & 6 \\
3 & 9 \\
4 & 12 \\
\bottomrule
\end{tabular}
\end{center}
\begin{enumerate}[label=(\alph*)]
  \item What is the acceleration?
  \item Plot the velocity-time graph. What is its shape?
  \item The area under the velocity-time graph gives the displacement. Estimate the displacement from $t=0$ to $t=4$.
\end{enumerate}

\problem An object has position $x(t) = 2t^3 - 3t^2 + t$.
\begin{enumerate}[label=(\alph*)]
  \item Find the velocity $v(t)$.
  \item Find the acceleration $a(t)$.
  \item At what time is the velocity zero?
  \item At what time is the acceleration zero?
\end{enumerate}
\end{explobox}

\keyidea{Position, velocity, and acceleration form a \emph{chain}: each is the rate of change of the previous one. Velocity is the slope of position; acceleration is the slope of velocity. This chain is the mathematical language of motion.}

% ────────────────────────────────────────────────────────────────────────────

\end{document}
